% ----------------------------------------------------------------------
% Tool lemma: value-density of bounded strategies (truncation).
% This file is an *artifact* for Section 4.3 of the thesis (tools for the
% repair). Under (A1) and (A6), the hedging functional pi equals the
% infimum of the OCE objective over bounded unconstrained strategies.
% This repairs failure point (F1): one may always choose a near-optimal
% target that is bounded (hence integrable) before invoking Hornik.
% When integrating into the final thesis, strip the standalone preamble
% and include this under the same numbering as the surrounding document
% (keep the "% draft: ..." comment to preserve the original label).
% Cross-references below are faked (hardcoded) so this artifact compiles
% without defining the targets; comments say what they should point to.
% ----------------------------------------------------------------------
\documentclass[11pt,a4paper]{report}
\usepackage[T1]{fontenc}
\usepackage[utf8]{inputenc}
\usepackage{amsmath,amssymb,amsthm}
\usepackage{enumitem}

\numberwithin{equation}{chapter}

\theoremstyle{plain}
\newtheorem{proposition}{Proposition}[chapter]
\newtheorem{lemma}[proposition]{Lemma}
\newtheorem{theorem}[proposition]{Theorem}
\newtheorem{corollary}[proposition]{Corollary}
\theoremstyle{definition}
\newtheorem{assumption}[proposition]{Assumption}
\newtheorem{definition}[proposition]{Definition}
\newtheorem{example}[proposition]{Example}
% Upright body, bold head (not the italic amsthm "remark" style).
\newtheorem{remark}[proposition]{Remark}

\newcommand{\E}{\mathbb{E}}
\newcommand{\PP}{\mathbb{P}}
\newcommand{\R}{\mathbb{R}}
\newcommand{\N}{\mathbb{N}}
\newcommand{\F}{\mathcal{F}}
\newcommand{\HH}{\mathcal{H}}
\newcommand{\Hu}{\mathcal{H}^{\mathrm{u}}}
\newcommand{\NNet}{\mathcal{N\!N}}
\newcommand{\as}{\ \PP\text{-a.s.}}

\begin{document}

\setcounter{chapter}{3}

\chapter{Artifact: value-density of bounded strategies}

\subsection*{Bounded strategies are value-dense under (A1) and (A6)}

% FAKE REF: "(F1)" -> the failure-point label of Section 4.2
%   (prop4_3_finite.tex): integrability of the representing maps f_k.
% FAKE REF: "(A1)", "(A6)" -> Chapter 3, Section 3.4 (full statements).
% FAKE REF: "Lemma~4.4" -> replace by \ref{lem:usc}
%   (lemma4_4_usc.tex): constant-envelope upper semicontinuity of OCEs.
On a general probability space, near-optimal strategies need not have
integrable hedge ratios, so Hornik's theorem cannot be applied
directly. This is failure point (F1). We repair it under two hypotheses
from Chapter~3: (A1) requires that the liability $Z$ is essentially
bounded, i.e.\ $Z\in L^\infty$, and (A6) requires that every admissible
gain is nonnegative, i.e.\ $W(\delta)\ge0$ a.s.\ for all $\delta\in\HH$.
Together they give $X+W(\eta)=-Z+W(\eta)\ge-\|Z\|_\infty$ for every
admissible $\eta$, which is the deterministic lower bound needed for
Lemma~4.4. The next lemma shows that, under (A1) and (A6), bounded
strategies are value-dense: every strategy can be replaced by a bounded
one without worsening the OCE objective by more than an arbitrary
$\varepsilon>0$. In particular, $\pi$ equals the infimum of the
objective over bounded strategies alone.

% draft: Lemma 4.6 (truncation / value-density of bounded strategies);
% same statement as lem:truncation-A in extension_with_A1_to_A6.tex.
\begin{lemma}\label{lem:truncation}
  Assume (A1) and (A6), and let $X=-Z$. Then for every $\delta\in\Hu$ and
  every $\varepsilon>0$ there exists a bounded $\delta'\in\Hu$ with
  \[
    \rho_\ell\bigl(X+W(H\circ\delta')\bigr)
    \le\rho_\ell\bigl(X+W(H\circ\delta)\bigr)+\varepsilon.
  \]
  In particular,
  \[
    \pi(X)=\inf\bigl\{\rho_\ell\bigl(X+W(H\circ\delta)\bigr):
    \delta\in\Hu,\ \delta\text{ bounded}\bigr\}.
  \]
\end{lemma}

\begin{proof}
  Fix $\delta\in\Hu$ and write $\eta:=H\circ\delta$, $Y:=X+W(\eta)$.
  Define the componentwise truncations $\delta^{(j)}=(\delta^{(j)}_k)_{k=0,\dots,n-1}$
  for $j\in\N$ by
  \[
    \delta^{(j)}_k:=
    \begin{cases}
      \delta_k & \text{if }\delta_k\in[-j,j], \\
      -j       & \text{if }\delta_k<-j,       \\
      j        & \text{if }\delta_k>j,
    \end{cases}
    \qquad k=0,\dots,n-1
  \]
  (applied componentwise if $d>1$). Then each $\delta^{(j)}$ is adapted
  (hence in $\Hu$) and bounded by $j$. Set
  $\eta^{(j)}:=H\circ\delta^{(j)}$ and $Y_j:=X+W(\eta^{(j)})$.

  Since there are finitely many time steps, there exists a (random) index
  $j_0<\infty$, namely
  \[
    j_0
    :=
    \max_{\substack{0\le k\le n-1\\ 1\le i\le d}}
    \bigl\lceil\bigl|\delta_k^i\bigr|\bigr\rceil,
  \]
  such that $\delta^{(j)}_k=\delta_k$ for all $k$ and all $j\ge j_0$. In
  particular $\delta^{(j)}_k\to\delta_k$ pointwise for every $k$. The
  recursive definition of $H\circ{}$ then yields $\eta^{(j)}\to\eta$
  pointwise, and therefore $Y_j\to Y$ pointwise.

  By (A6), $W(\eta^{(j)})\ge0$ a.s., so $Y_j\ge-\|Z\|_\infty$ a.s.\ by
  (A1). Lemma~4.4 with $B=\|Z\|_\infty$ gives
  \[
    \limsup_{j\to\infty}\rho_\ell(Y_j)\le\rho_\ell(Y).
  \]
  Choose $j$ large enough that $\rho_\ell(Y_j)\le\rho_\ell(Y)+\varepsilon$,
  and set $\delta':=\delta^{(j)}$.
\end{proof}

\end{document}
